Les suites introduisent la limite, et donc le premier vrai changement de nature de l'analyse.

Le chapitre commence par vérifier que la convergence de la bibliothèque — écrite avec des
filtres, pour valoir dans n'importe quel espace — est bien la définition avec $\varepsilon$ du
programme, et que la divergence vers $+\infty$ est bien celle avec $A$. Sans cette
vérification, tout ce qui suit porterait sur un objet dont rien ne garantit qu'il est celui
du cours. Ce premier énoncé n'apprend rien de neuf ; il rend les autres lisibles.

Vient ensuite ce que le programme admet et qui est ici démontré : toute suite croissante
majorée converge. C'est la propriété de la borne supérieure de $\mathbb{R}$, c'est-à-dire sa
complétude, qui n'est jamais énoncée au lycée.

Un énoncé reste incomplet : les suites adjacentes sont définies et leurs premières propriétés
établies, mais la dichotomie — l'usage qu'on en fait pour encadrer une solution — n'est pas
traitée.
